\subsection{Pregunta 4}

¿En qué escenarios es bueno compartir nuestra información públicamente? ¿O en cuáles no? Justifica.

\begin{respuesta}
    Depende de para qué lo hagas y de qué tanto te pueda costar. Compartir en público
    está padre cuando te conviene que te encuentren: subir tu perfil en LinkedIn o tu
    GitHub para que te contraten, dar a conocer un proyecto o un negocio, divulgar lo
    que sabes o apoyar una causa. Ahí lo que publicas es justo lo que quieres que se
    vea, y el riesgo es bajo.

    En cambio, no conviene cuando lo que subes ayuda a que alguien te ubique o se haga
    pasar por ti. Por ejemplo: dónde vives, tus rutinas, cuándo sales de viaje (y más si
    lo subes en tiempo real), fotos de tus documentos o boletos, o datos de menores. Un
    dato suelto parece inofensivo, pero al juntar varios, como hicimos en esta práctica,
    se arma un perfil completísimo. La regla sencilla es: si no te gustaría que un
    desconocido lo leyera, mejor no lo publiques.
\end{respuesta}
